% (C) 2026 Ralf Stephan, in collaboration with Claude Code.
% Released under CC0 1.0 Universal (public-domain dedication).
% See https://creativecommons.org/publicdomain/zero/1.0/

\documentclass[11pt]{article}
\usepackage[T1]{fontenc}
\usepackage[utf8]{inputenc}
\usepackage{lmodern}
\usepackage[margin=2.5cm]{geometry}
\usepackage{amsmath,amssymb}
\usepackage{braket}
\usepackage[hidelinks]{hyperref}

% Lean names and paths: verbatim, breakable at . _ /
\DeclareUrlCommand\lean{\urlstyle{tt}}
\def\UrlBreaks{\do\.\do\_\do\/}
% Marks an open problem that blocks checking a main claim of the brief
\newcommand{\blocker}[1]{\textsc{[blocks verification: #1]}~}
% Extra stretch for paragraphs with long, unbreakable Lean names
\emergencystretch=6em

\title{Open problems:\\ Non-linear representations in quantum mechanics}
\author{}
\date{}

\begin{document}
\maketitle

\noindent
Open problems raised by the technical brief in this directory, and by its Lean formalization
(see \lean{formalization-log.tex}).
Items marked \textsc{[blocks verification]} stand in the way of checking a main
claim of the brief as the brief states it: a \emph{formalization gap} can be closed
in Lean alone, while a \emph{brief issue} needs the brief to be corrected or made
precise first.

\section*{From the brief}
\begin{enumerate}
\item \textbf{Madelung representation.} \emph{Research.} ``The Madelung equations \ldots\
  provide such a representation'', one ``not even tied to vector spaces''. The equivalence of
  the Madelung equations with quantum theory, and in what sense they represent QST:H, are not
  shown. As stated the equivalence is false: the Madelung equations have solutions that come
  from no wave function, and to exclude them one must add a quantization condition on the
  circulation of the velocity field, $\oint \nabla S \cdot d\ell \in 2\pi\hbar\,\mathbb{Z}$,
  which is not part of the hydrodynamic equations and is hard to state at the nodes, where the
  density vanishes and the phase $S$ is undefined (T.~C.~Wallstrom, Phys.\ Rev.\ A \textbf{49}
  (1994), 1613). A precise statement is needed first, and checking it needs partial
  differential equations on $\mathbb{R}^3$ with a singular nodal set, which neither Mathlib nor
  TauCeti provides. (Brief, Sec.~5.)
\end{enumerate}

\section*{From the formalization}
\begin{enumerate}
\item None found. The blocker of the previous version of the brief (Theorem~1 as an
  independence statement) is resolved: the revised brief states its mathematical content as
  Proposition~1 and Corollary~3, which are checked in Lean, and argues Theorem~4 from them
  informally, which is checked by reading (Log, ``Differences from the brief'').
\end{enumerate}

\end{document}
